\chapter{تخصيص السنة لعموم القرآن ووجوب طاعة الرسول}

\section{مقدمة في تخصيص السنة للقرآن}
الحمد لله والصلاة والسلام على رسول الله، أما بعد؛ فقد تناولنا في الدرس السابق ما نزل من القرآن عاماً وأقسامه. ونشرع في هذا الدرس في بيان قاعدة أصولية جليلة قررها الإمام الشافعي رحمه الله، وهي: «ما نزل من القرآن عاماً ودلت السنة على أنه يراد به الخاص». فالسنة المطهرة تأتي لتخصص عموم القرآن الكريم وتُقيد مطلقه، وقد ضرب الشافعي لذلك ستة أمثلة فقهية تطبيقية:

\section{الأمثلة التطبيقية على تخصيص السنة للقرآن}

\subsection{المثال الأول: آيات المواريث والوصية}
قال تعالى: ﴿وَلِأَبَوَيْهِ لِكُلِّ وَاحِدٍ مِّنْهُمَا السُّدُسُ مِمَّا تَرَكَ إِن كَانَ لَهُ وَلَدٌ﴾، وقال: ﴿وَلَكُمْ نِصْفُ مَا تَرَكَ أَزْوَاجُكُمْ﴾. هذه الآيات جاءت بلفظ عام يشمل كل والد وولد وكل زوج وزوجة. لكن السنة المطهرة خصصت هذا العموم، فبينت أنه لا يتوارث أهل ملتين (المسلم والكافر)، وأن القاتل لا يرث.

كذلك في الوصية، قال تعالى: ﴿مِن بَعْدِ وَصِيَّةٍ يُوصَىٰ بِهَا أَوْ دَيْنٍ﴾، فالآية عامة في أي مقدار للوصية، لكن السنة خصصتها بحديث سعد بن أبي وقاص رضي الله عنه: «الثلث والثلث كثير»، فلا تنفذ الوصية فيما زاد عن الثلث. كما بينت السنة وأجمع المسلمون على أن سداد الدين مُقَدَّم على الوصية، رغم أن الآية ذكرت الوصية أولاً.

\subsection{المثال الثاني: آية الوضوء وغسل القدمين}
قال تعالى: ﴿إِذَا قُمْتُمْ إِلَى الصَّلَاةِ فَاغْسِلُوا وُجُوهَكُمْ وَأَيْدِيَكُمْ إِلَى الْمَرَافِقِ وَامْسَحُوا بِرُءُوسِكُمْ وَأَرْجُلَكُمْ إِلَى الْكَعْبَيْنِ﴾. قصد الله عز وجل غسل القدمين كغسل الوجه واليدين، وظاهر الآية أنه لا يجزئ إلا الغسل لجميع المتوضئين. لكن سنة النبي صلى الله عليه وسلم بتواتر مسحه على الخفين، دلت على أن المراد بغسل القدمين بعض المتوضئين دون بعض، فمن لبس الخفين على طهارة خُصَّ بجواز المسح، ومن لم يلبسهما وجب عليه الغسل.

\subsection{المثال الثالث: آية حد السرقة}
قال تعالى: ﴿وَالسَّارِقُ وَالسَّارِقَةُ فَاقْطَعُوا أَيْدِيَهُمَا﴾. لفظ «السارق» عام، ولفظ «اليد» عام (يشمل من أطراف الأصابع إلى الإبط). فجاءت السنة المطهرة لتخصص السارق بمن سرق مالاً محرزاً يبلغ ربع دينار فصاعداً (فلا قطع في ثمر معلق ولا في أقل من ربع دينار)، وخصصت السنة اليد بـ «الكف» ليُقطع من المفصل (الرسغ) وليس من المرفق أو العضد.

\subsection{المثال الرابع: آية حد الزنا}
قال تعالى: ﴿الزَّانِيَةُ وَالزَّانِي فَاجْلِدُوا كُلَّ وَاحِدٍ مِّنْهُمَا مِائَةَ جَلْدَةٍ﴾، وهذا عام. لكن السنة دلت على أن المراد بجلد المائة هما «الحران البكران» فقط، لأن النبي صلى الله عليه وسلم رجم الثيب المحصن ولم يجلده (كما في قصة ماعز والغامدية). أما الإماء والعبيد، فقد بين القرآن أن عليهم نصف ما على المحصنات (أي خمسين جلدة)، فخُصصت الآية الكريمة بالسنة.

\subsection{المثال الخامس: سهم ذوي القربى في الغنائم}
قال تعالى: ﴿وَاعْلَمُوا أَنَّمَا غَنِمْتُم مِّن شَيْءٍ فَأَنَّ لِلَّهِ خُمُسَهُ وَلِلرَّسُولِ وَلِذِي الْقُرْبَىٰ﴾. «ذو القربى» لفظ عام قد يشمل كل قريش، أو كل بني عبد مناف. لكن السنة خصصته بطائفتين فقط هما: بنو هاشم وبنو المطلب. ولما سأل عثمان بن عفان وجبير بن مطعم (وهما من بني عبد شمس وبني نوفل أبناء عبد مناف) عن سبب إعطاء بني المطلب ومنعهم مع أن قرابتهم واحدة، شبّك النبي صلى الله عليه وسلم أصابعه وقال: «إنما بنو هاشم وبنو المطلب شيء واحد في الجاهلية والإسلام». فدلت السنة على تخصيص ذوي القربى بهاتين البطنتين.

\subsection{المثال السادس: استثناء السَّلَب من عموم الغنائم}
الآية السابقة جعلت الغنائم مقسومة (يُؤخذ خمسها ويوزع الباقي)، لكن السنة خصت «السَّلَب» (وهو ما يكون على المقتول من سلاح وعتاد ومال) من عموم الغنيمة، فجعلته للقاتل خالصاً إذا قتل مشركاً في إقبال وحالة قتال. ودليل ذلك قضاء النبي صلى الله عليه وسلم لـ أبي قتادة رضي الله عنه يوم حنين بسلب المشرك الذي قتله بعد أن أقام البينة، وقال: «من قتل قتيلاً له عليه بينة فله سلبه». فهذا السلب لا يُخمّس ولا يدخل في عموم الغنيمة.

\section{فرضية طاعة رسول الله صلى الله عليه وسلم}
بعد أن بيّن الشافعي أمثلة تخصيص السنة للقرآن، أراد أن يُؤصل مسألة في غاية الأهمية، وهي أن اتباع النبي صلى الله عليه وسلم فرض حتمي، وأن سنته مستقلة بالتشريع.

\subsection{اقتران الإيمان بالله بالإيمان بالرسول}
جعل الله الإيمان برسوله صلى الله عليه وسلم كمالاً وشرطاً للإيمان به جل وعلا. قال تعالى: ﴿إِنَّمَا الْمُؤْمِنُونَ الَّذِينَ آمَنُوا بِاللَّهِ وَرَسُولِهِ﴾، فلو آمن عبد بالله ولم يؤمن برسوله لم يقع عليه اسم الإيمان أبداً.

\textbf{تنبيه منهجي:} أورد الإمام الشافعي رحمه الله في هذا الموضع استدلالاً بقوله تعالى: ﴿فَآمِنُوا بِاللَّهِ وَرَسُولِهِ﴾، مع أن الآية في سورة النساء جاءت بلفظ الجمع: ﴿فَآمِنُوا بِاللَّهِ وَرُسُلِهِ﴾، والآيات الدالة على الإفراد كثيرة في سور أخرى (كالأعراف والتغابن). وقد أشار العلامة أحمد شاكر رحمه الله إلى أن هذا وهم وسبق لسان من الإمام الشافعي، وأن بقاءه في المخطوطات قروناً ينسخه العلماء دون تنبه يدل على خطورة «التقليد» حتى لكبار الأئمة. وهذا تطبيق عملي لقول الشافعي نفسه: «وبالتقليد أغفل من أغفل منهم، والله يغفر لنا ولهم»، فالعصمة للوحيين فقط.

\subsection{السنة هي «الحكمة» في القرآن}
قال تعالى: ﴿لَقَدْ مَنَّ اللَّهُ عَلَى الْمُؤْمِنِينَ إِذْ بَعَثَ فِيهِمْ رَسُولًا مِّنْ أَنفُسِهِمْ يَتْلُو عَلَيْهِمْ آيَاتِهِ وَيُزَكِّيهِمْ وَيُعَلِّمُمُ الْكِتَابَ وَالْحِكْمَةَ﴾، وقال موجهاً الخطاب لنساء النبي: ﴿وَاذْكُرْنَ مَا يُتْلَىٰ فِي بُيُوتِكُنَّ مِنْ آيَاتِ اللَّهِ وَالْحِكْمَةِ﴾.
قال الشافعي: «سمعت من أرضى من أهل العلم بالقرآن يقول: الحكمة سُنَّة رسول الله صلى الله عليه وسلم». فإن الله منّ على خلقه بتعليمهم الكتاب (وهو القرآن) والحكمة (وهي السنة)، ولا يجوز أن يقال الحكمة هنا إلا سنة رسول الله صلى الله عليه وسلم لأنها قُرنت بكتاب الله.

\subsection{الرد على منكري السنة (القرآنيين)}
إن من يدعي الاكتفاء بالقرآن ليرد سنة النبي صلى الله عليه وسلم هو في الحقيقة مكذب للقرآن. فالقرآن هو الذي قال: ﴿وَمَا آتَاكُمُ الرَّسُولُ فَخُذُوهُ وَمَا نَهَاكُمْ عَنْهُ فَانتَهُوا﴾.
فمن قَبِلَ عن رسول الله صلى الله عليه وسلم سنته، فقد قَبِلَ عن الله فَرْضَه. ومن فرَّق بينهما وقَبِلَ واحداً وترك الآخر، فقد ردَّ على الله أمره. وسنة رسول الله صلى الله عليه وسلم مُبَيِّنةٌ عن اللهِ مرادَه، دليلةٌ على خاصِّه وعامِّه، ولا يجوز لأحد أن يُخالفها كائناً من كان.